% =========================================================
% Terms and Formulas
% =========================================================

\section{Syntax of First-Order Logic: Terms and Formulas}

\begin{remark*}[Section toolkit: Syntax ledger]
This section fixes the syntactic inventory for first-order logic before any
semantic machinery is introduced. The ledger proceeds from symbols to terms,
from terms to atomic formulas, and from atomic formulas to all well-formed
formulas.

\begin{center}
\small
\begin{tabular}{l l l}
\toprule
\textbf{Layer} & \textbf{Item} & \textbf{Role} \\
\midrule
Symbol stock & \hyperref[def:variable]{Variable} & reusable placeholder symbol \\
Symbol stock & \hyperref[def:constant-symbol]{Constant Symbol} & nullary name-forming symbol \\
Symbol stock & \hyperref[def:function-symbol]{Function Symbol} & operation-forming symbol \\
Symbol stock & \hyperref[def:predicate-symbol]{Predicate Symbol} & assertion-forming symbol \\
Object syntax & \hyperref[def:term]{Term} & expression that names an object syntactically \\
Formula syntax & \hyperref[def:atomic-formula]{Atomic Formula} & predicate symbol applied to terms \\
Formula syntax & \hyperref[def:wff-pred]{Well-Formed Formula} & formula generated by the grammar \\
Formula syntax & \hyperref[def:molecular]{Molecular Formula} & non-atomic well-formed formula \\
\bottomrule
\end{tabular}
\end{center}
\end{remark*}

\begin{remark*}[Exposition]
The grammar has two tracks. Terms are built first, because atomic formulas need
terms as their inputs. Formulas are built second, because they begin with atomic
formulas and then close under logical connectives and quantifiers. Semantic
notions such as structures, assignments, satisfaction, truth, and models enter
later; the present section records only which strings count as grammatical
expressions.
\end{remark*}

\begin{remark*}[Grammar boundary]
Terms name objects. Formulas make assertions. A term is not a formula, and a
formula is not a term. For example, \(f(x)\) is a term, while \(P(f(x))\) is a
formula.
\end{remark*}

\begin{remark*}[Predicate-symbol convention]
Object-language predicate symbols, such as \(P,Q,R\), are symbols that occur
inside first-order formulas. Project metalanguage predicate readings, such as
\(\operatorname{Term}(t)\) or \(\operatorname{AtomicFormula}(\varphi)\), are
external classification predicates used by these notes to describe the grammar.
The two levels should not be confused.
\end{remark*}

% ---------------------------------------------------------
% Variable
% ---------------------------------------------------------
\begin{definitionbox}{Definition (Variable)}
\begin{definition}[Variable]
\label{def:variable}
For a first-order language \(\mathcal L\), a \emph{variable} is a symbol drawn
from the distinguished variable stock of \(\mathcal L\). Variables are usually
written
\[
x,y,z,x_0,x_1,x_2,\ldots .
\]
\end{definition}
\end{definitionbox}

\begin{remark*}[Standard quantified statement]
\[
\mathsf{Var}_{\mathcal L}=\{x,y,z,x_0,x_1,x_2,\ldots\}.
\]
\end{remark*}

\begin{remark*}[Definition predicate reading]
\(\operatorname{Variable}(x)\) means \(x\in\mathsf{Var}_{\mathcal L}\).
\end{remark*}

\begin{remark*}[Negated quantified statement]
\[
x\notin\mathsf{Var}_{\mathcal L}.
\]
\end{remark*}

\begin{remark*}[Negation predicate reading]
\(\neg\operatorname{Variable}(x)\) means that \(x\) is not one of the variable
symbols of \(\mathcal L\).
\end{remark*}

\begin{remark*}[Failure modes]
A symbol fails to be a variable when it belongs to a different syntactic role,
such as a constant symbol, function symbol, predicate symbol, connective, or
punctuation mark.
\end{remark*}

\begin{remark*}[Failure mode decomposition]
\[
\underbrace{c}_{\text{constant symbol}}
\qquad
\underbrace{f}_{\text{function symbol}}
\qquad
\underbrace{P}_{\text{predicate symbol}}
\qquad
\underbrace{\forall}_{\text{logical symbol}}.
\]
\end{remark*}

\begin{remark*}[Interpretation]
A variable is a basic syntactic placeholder. It can occur inside terms and
formulas, and later sections explain how quantifiers bind occurrences of
variables.
\end{remark*}

\begin{remark*}[Examples]
\(x\), \(y\), \(z\), \(x_0\), and \(x_1\) are standard variable symbols.
\end{remark*}

\begin{remark*}[Non-Examples]
\(c\), \(0\), \(f\), \(P\), \(\forall\), and \(\land\) are not variables merely
because they are symbols of the language.
\end{remark*}

\begin{dependencies}
\begin{itemize}
  \item \hyperref[def:language]{Language}
\end{itemize}
\end{dependencies}

% ---------------------------------------------------------
% Constant Symbol
% ---------------------------------------------------------
\begin{definitionbox}{Definition (Constant Symbol)}
\begin{definition}[Constant Symbol]
\label{def:constant-symbol}
For a first-order language \(\mathcal L\), a \emph{constant symbol} is a
non-logical symbol of arity \(0\). A constant symbol may occur as a term without
taking input terms.
\end{definition}
\end{definitionbox}

\begin{remark*}[Standard quantified statement]
\[
\mathsf{Const}_{\mathcal L}
=
\{c\in\mathcal L : \operatorname{arity}(c)=0\}.
\]
\end{remark*}

\begin{remark*}[Definition predicate reading]
\(\operatorname{ConstantSymbol}(c)\) means
\(c\in\mathsf{Const}_{\mathcal L}\).
\end{remark*}

\begin{remark*}[Negated quantified statement]
\[
c\notin\mathsf{Const}_{\mathcal L}.
\]
\end{remark*}

\begin{remark*}[Negation predicate reading]
\(\neg\operatorname{ConstantSymbol}(c)\) means that \(c\) is not a nullary
non-logical symbol of \(\mathcal L\).
\end{remark*}

\begin{remark*}[Failure modes]
A symbol fails to be a constant symbol when it is not in the language, is
logical punctuation, or has positive arity.
\end{remark*}

\begin{remark*}[Failure mode decomposition]
\[
\underbrace{x}_{\text{variable}}
\qquad
\underbrace{f}_{\text{positive arity}}
\qquad
\underbrace{(}_{\text{punctuation}}.
\]
\end{remark*}

\begin{remark*}[Interpretation]
A constant symbol is a simple name-forming symbol in the grammar. It supplies a
base case for terms alongside variables.
\end{remark*}

\begin{remark*}[Examples]
In an arithmetic language, \(0\) and \(1\) are commonly used as constant
symbols.
\end{remark*}

\begin{remark*}[Non-Examples]
\(f(x)\) is not a constant symbol; it is a term built from a function symbol
and an input term.
\end{remark*}

\begin{dependencies}
\begin{itemize}
  \item \hyperref[def:variable]{Variable}
\end{itemize}
\end{dependencies}

% ---------------------------------------------------------
% Function Symbol
% ---------------------------------------------------------
\begin{definitionbox}{Definition (Function Symbol)}
\begin{definition}[Function Symbol]
\label{def:function-symbol}
For a first-order language \(\mathcal L\), an \(n\)-ary \emph{function symbol}
is a non-logical symbol of arity \(n\geq 1\). If \(f\) has arity \(n\), then it
forms terms by the pattern
\[
f(t_1,\ldots,t_n).
\]
\end{definition}
\end{definitionbox}

\begin{remark*}[Standard quantified statement]
\[
\mathsf{Func}_{\mathcal L}^{(n)}
=
\{f\in\mathcal L : \operatorname{arity}(f)=n\text{ and }n\geq 1\}.
\]
\end{remark*}

\begin{remark*}[Definition predicate reading]
\(\operatorname{FunctionSymbol}(f)\) means
\(f\in\mathsf{Func}_{\mathcal L}^{(n)}\) for some \(n\geq 1\).
\end{remark*}

\begin{remark*}[Negated quantified statement]
\[
\forall n\geq 1\;(f\notin\mathsf{Func}_{\mathcal L}^{(n)}).
\]
\end{remark*}

\begin{remark*}[Negation predicate reading]
\(\neg\operatorname{FunctionSymbol}(f)\) means that \(f\) is not a positive-arity
function symbol of \(\mathcal L\).
\end{remark*}

\begin{remark*}[Failure modes]
A symbol fails to be a function symbol when it is a variable, constant symbol,
predicate symbol, logical symbol, punctuation mark, or has no assigned positive
arity.
\end{remark*}

\begin{remark*}[Failure mode decomposition]
\[
\underbrace{c}_{\text{arity }0}
\qquad
\underbrace{P}_{\text{predicate symbol}}
\qquad
\underbrace{\neg}_{\text{logical symbol}}.
\]
\end{remark*}

\begin{remark*}[Interpretation]
A function symbol is an operation-forming symbol in the grammar. It is not a
term by itself; it becomes part of a term only after it is supplied the required
number of input terms.
\end{remark*}

\begin{remark*}[Examples]
In an arithmetic language, \(S\) may be unary, while \(+\) and \(\cdot\) may be
binary function symbols.
\end{remark*}

\begin{remark*}[Non-Examples]
\(f(x)\) is not a function symbol. It is a term formed by applying the function
symbol \(f\) to the term \(x\).
\end{remark*}

\begin{dependencies}
\begin{itemize}
  \item \hyperref[def:constant-symbol]{Constant Symbol}
\end{itemize}
\end{dependencies}

% ---------------------------------------------------------
% Predicate Symbol
% ---------------------------------------------------------
\begin{definitionbox}{Definition (Predicate Symbol)}
\begin{definition}[Predicate Symbol]
\label{def:predicate-symbol}
For a first-order language \(\mathcal L\), an \(n\)-ary \emph{predicate symbol}
is a non-logical symbol of arity \(n\geq 1\). If \(P\) has arity \(n\), then it
forms atomic formulas by the pattern
\[
P(t_1,\ldots,t_n).
\]
\end{definition}
\end{definitionbox}

\begin{remark*}[Standard quantified statement]
\[
\mathsf{Pred}_{\mathcal L}^{(n)}
=
\{P\in\mathcal L : \operatorname{arity}(P)=n\text{ and }n\geq 1\}.
\]
\end{remark*}

\begin{remark*}[Definition predicate reading]
\(\operatorname{PredicateSymbol}(P)\) means
\(P\in\mathsf{Pred}_{\mathcal L}^{(n)}\) for some \(n\geq 1\).
\end{remark*}

\begin{remark*}[Negated quantified statement]
\[
\forall n\geq 1\;(P\notin\mathsf{Pred}_{\mathcal L}^{(n)}).
\]
\end{remark*}

\begin{remark*}[Negation predicate reading]
\(\neg\operatorname{PredicateSymbol}(P)\) means that \(P\) is not a positive-arity
predicate symbol of \(\mathcal L\).
\end{remark*}

\begin{remark*}[Failure modes]
A symbol fails to be a predicate symbol when it belongs to the term-forming
side of the language, is a logical symbol, is punctuation, or has no assigned
positive arity as a predicate symbol.
\end{remark*}

\begin{remark*}[Failure mode decomposition]
\[
\underbrace{x}_{\text{variable}}
\qquad
\underbrace{c}_{\text{constant symbol}}
\qquad
\underbrace{f}_{\text{function symbol}}
\qquad
\underbrace{\exists}_{\text{logical symbol}}.
\]
\end{remark*}

\begin{remark*}[Interpretation]
Predicate symbols are the atomic-formula-forming vocabulary of the language.
They are syntactic symbols, not project metalanguage predicates.
\end{remark*}

\begin{remark*}[Examples]
\(P\), \(Q\), \(R\), \(<\), and \(=\) may be predicate symbols when the language
declares the corresponding arities.
\end{remark*}

\begin{remark*}[Non-Examples]
\(P(x)\) is not a predicate symbol. It is an atomic formula formed from the
predicate symbol \(P\) and the term \(x\).
\end{remark*}

\begin{dependencies}
\begin{itemize}
  \item \hyperref[def:function-symbol]{Function Symbol}
\end{itemize}
\end{dependencies}

% ---------------------------------------------------------
% Term
% ---------------------------------------------------------
\begin{definitionbox}{Definition (Term)}
\begin{definition}[Term]
\label{def:term}
The set of \emph{terms} of a first-order language \(\mathcal L\), denoted
\(\mathsf{Term}_{\mathcal L}\), is the smallest set of strings satisfying:
\begin{enumerate}
\item every variable is a term;
\item every constant symbol is a term;
\item if \(f\) is an \(n\)-ary function symbol and
      \(t_1,\ldots,t_n\) are terms, then \(f(t_1,\ldots,t_n)\) is a term;
\item no string is a term unless it is generated by finitely many applications
      of the previous clauses.
\end{enumerate}
\end{definition}
\end{definitionbox}

\begin{remark*}[Standard quantified statement]
\[
\mathsf{Term}_{\mathcal L}
=
\bigcap
\left\{
S:
\mathsf{Var}_{\mathcal L}\subseteq S,\;
\mathsf{Const}_{\mathcal L}\subseteq S,\;
(\forall n\geq 1)(\forall f\in\mathsf{Func}_{\mathcal L}^{(n)})
(\forall t_1,\ldots,t_n\in S)
f(t_1,\ldots,t_n)\in S
\right\}.
\]
\end{remark*}

\begin{remark*}[Definition predicate reading]
\(\operatorname{Term}(t)\) means \(t\in\mathsf{Term}_{\mathcal L}\).
\end{remark*}

\begin{remark*}[Negated quantified statement]
\[
\sigma\notin\mathsf{Term}_{\mathcal L}
\Longleftrightarrow
\exists S
\left[
\begin{aligned}
&\mathsf{Var}_{\mathcal L}\subseteq S
\land
\mathsf{Const}_{\mathcal L}\subseteq S\\
&\land
(\forall n\geq 1)(\forall f\in\mathsf{Func}_{\mathcal L}^{(n)})
(\forall t_1,\ldots,t_n\in S)
f(t_1,\ldots,t_n)\in S\\
&\land\ \sigma\notin S
\end{aligned}
\right].
\]
\end{remark*}

\begin{remark*}[Negation predicate reading]
\(\neg\operatorname{Term}(\sigma)\) means that \(\sigma\) is not generated by
the term-formation clauses.
\end{remark*}

\begin{remark*}[Failure modes]
A string fails to be a term when it uses a predicate symbol, connective, or
quantifier at the outer syntactic level, when a function symbol has too few or
too many input terms, or when punctuation does not match the arity pattern.
\end{remark*}

\begin{remark*}[Failure mode decomposition]
\[
\underbrace{P(x)}_{\text{atomic formula, not a term}}
\qquad
\underbrace{f(x,y)}_{\text{wrong arity if }f\text{ is unary}}
\qquad
\underbrace{\forall x}_{\text{quantifier fragment}}.
\]
\end{remark*}

\begin{remark*}[Interpretation]
Terms are the noun-like expressions of first-order syntax. They provide the
inputs to predicate symbols but do not themselves make assertions.
\end{remark*}

\begin{remark*}[Examples]
\(x\), \(c\), \(f(x)\), \(g(c,x)\), and \(h(f(x),c)\) are terms when the
language assigns the displayed arities.
\end{remark*}

\begin{remark*}[Non-Examples]
\(P(x)\), \(\neg P(x)\), \(\forall x\,P(x)\), and \(x\land y\) are not terms.
\end{remark*}

\begin{dependencies}
\begin{itemize}
  \item \hyperref[def:variable]{Variable}
  \item \hyperref[def:constant-symbol]{Constant Symbol}
  \item \hyperref[def:function-symbol]{Function Symbol}
\end{itemize}
\end{dependencies}

% ---------------------------------------------------------
% Atomic Formula
% ---------------------------------------------------------
\begin{definitionbox}{Definition (Atomic Formula)}
\begin{definition}[Atomic Formula]
\label{def:atomic-formula}
An \emph{atomic formula} of a first-order language \(\mathcal L\) is a string
of the form
\[
P(t_1,\ldots,t_n),
\]
where \(P\) is an \(n\)-ary predicate symbol and
\(t_1,\ldots,t_n\) are terms.
\end{definition}
\end{definitionbox}

\begin{remark*}[Standard quantified statement]
\[
\mathsf{Atom}_{\mathcal L}
=
\left\{
P(t_1,\ldots,t_n):
P\in\mathsf{Pred}_{\mathcal L}^{(n)}
\text{ and }
t_1,\ldots,t_n\in\mathsf{Term}_{\mathcal L}
\right\}.
\]
\end{remark*}

\begin{remark*}[Definition predicate reading]
\(\operatorname{AtomicFormula}(\varphi)\) means
\(\varphi\in\mathsf{Atom}_{\mathcal L}\).
\end{remark*}

\begin{remark*}[Negated quantified statement]
\[
\varphi\notin\mathsf{Atom}_{\mathcal L}.
\]
\end{remark*}

\begin{remark*}[Negation predicate reading]
\(\neg\operatorname{AtomicFormula}(\varphi)\) means that \(\varphi\) is not a
predicate symbol applied to the required number of terms.
\end{remark*}

\begin{remark*}[Failure modes]
A string fails to be atomic when it does not begin with a predicate symbol,
when the predicate symbol receives the wrong number of term inputs, when an
input is not a term, or when a logical connective or quantifier has already
been applied.
\end{remark*}

\begin{remark*}[Failure mode decomposition]
\[
\underbrace{f(x)}_{\text{term, not formula}}
\qquad
\underbrace{P(x,y)}_{\text{wrong arity if }P\text{ is unary}}
\qquad
\underbrace{\neg P(x)}_{\text{molecular formula}}.
\]
\end{remark*}

\begin{remark*}[Interpretation]
Atomic formulas are the base cases of formula formation. They are the first
strings in the grammar that count as formulas.
\end{remark*}

\begin{remark*}[Examples]
\(P(x)\), \(R(x,c)\), \(x=y\), and \(x<y\) are atomic formulas when the
language supplies the relevant predicate symbols and arities.
\end{remark*}

\begin{remark*}[Non-Examples]
\(x\), \(f(x)\), \(\neg P(x)\), and \((P(x)\land Q(x))\) are not atomic
formulas.
\end{remark*}

\begin{dependencies}
\begin{itemize}
  \item \hyperref[def:predicate-symbol]{Predicate Symbol}
  \item \hyperref[def:term]{Term}
\end{itemize}
\end{dependencies}

% ---------------------------------------------------------
% Well-Formed Formula
% ---------------------------------------------------------
\begin{definitionbox}{Definition (Well-Formed Formula)}
\begin{definition}[Well-Formed Formula]
\label{def:wff-pred}
The set of \emph{well-formed formulas} of a first-order language
\(\mathcal L\), denoted \(\mathsf{Form}_{\mathcal L}\), is the smallest set of
strings satisfying:
\begin{enumerate}
\item every atomic formula is a well-formed formula;
\item if \(\varphi\) is a well-formed formula, then
      \(\neg\varphi\) is a well-formed formula;
\item if \(\varphi\) and \(\psi\) are well-formed formulas and
      \(\circ\in\{\land,\lor,\to,\leftrightarrow\}\), then
      \((\varphi\circ\psi)\) is a well-formed formula;
\item if \(\varphi\) is a well-formed formula and \(x\) is a variable, then
      \(\forall x\,\varphi\) and \(\exists x\,\varphi\) are well-formed
      formulas;
\item no string is a well-formed formula unless it is generated by finitely
      many applications of the previous clauses.
\end{enumerate}
\end{definition}
\end{definitionbox}

\begin{remark*}[Standard quantified statement]
\[
\begin{aligned}
\mathsf{Form}_{\mathcal L}
=
\bigcap\{S:\;&
\mathsf{Atom}_{\mathcal L}\subseteq S,\\
&(\forall\varphi\in S)(\neg\varphi\in S),\\
&(\forall\varphi,\psi\in S)
(\forall\circ\in\{\land,\lor,\to,\leftrightarrow\})
((\varphi\circ\psi)\in S),\\
&(\forall x\in\mathsf{Var}_{\mathcal L})(\forall\varphi\in S)
(\forall x\,\varphi\in S\land \exists x\,\varphi\in S)\}.
\end{aligned}
\]
\end{remark*}

\begin{remark*}[Definition predicate reading]
\(\operatorname{WellFormedFormula}(\varphi)\) means
\(\varphi\in\mathsf{Form}_{\mathcal L}\).
\end{remark*}

\begin{remark*}[Negated quantified statement]
\[
\varphi\notin\mathsf{Form}_{\mathcal L}
\Longleftrightarrow
\exists S\left[
\begin{aligned}
&\mathsf{Atom}_{\mathcal L}\subseteq S
\land
(\forall\psi\in S)(\neg\psi\in S)\\
&\land
(\forall\psi,\chi\in S)
(\forall\circ\in\{\land,\lor,\to,\leftrightarrow\})
((\psi\circ\chi)\in S)\\
&\land
(\forall x\in\mathsf{Var}_{\mathcal L})(\forall\psi\in S)
(\forall x\,\psi\in S\land \exists x\,\psi\in S)\\
&\land\ \varphi\notin S
\end{aligned}
\right].
\]
\end{remark*}

\begin{remark*}[Negation predicate reading]
\(\neg\operatorname{WellFormedFormula}(\varphi)\) means that \(\varphi\) is not
generated by the formula-formation clauses.
\end{remark*}

\begin{remark*}[Failure modes]
A string fails to be well formed when it is a bare term, has a missing operand,
uses a quantifier without a variable and formula, has unmatched punctuation, or
cannot be produced by finitely many formation steps.
\end{remark*}

\begin{remark*}[Failure mode decomposition]
\[
\underbrace{f(x)}_{\text{term only}}
\qquad
\underbrace{P(x)\land}_{\text{missing right formula}}
\qquad
\underbrace{\forall P(x)}_{\text{missing bound variable}}
\qquad
\underbrace{(P(x)}_{\text{unmatched punctuation}}.
\]
\end{remark*}

\begin{remark*}[Interpretation]
Well-formed formulas are exactly the grammatical formula strings of the
language. Later semantic and proof-theoretic sections apply only to these
strings.
\end{remark*}

\begin{remark*}[Examples]
\(P(x)\), \(\neg P(x)\), \((P(x)\to Q(x))\), and
\(\forall x\,(P(x)\to Q(x))\) are well-formed formulas.
\end{remark*}

\begin{remark*}[Non-Examples]
\(x\), \(f(x)\), \(P(x)\land\), and \(\forall P(x)\) are not well-formed
formulas.
\end{remark*}

\begin{dependencies}
\begin{itemize}
  \item \hyperref[def:variable]{Variable}
  \item \hyperref[def:atomic-formula]{Atomic Formula}
  \item \hyperref[def:universal-q]{Universal Formula}
  \item \hyperref[def:existential-q]{Existential Formula}
\end{itemize}
\end{dependencies}

% ---------------------------------------------------------
% Molecular Formula
% ---------------------------------------------------------
\begin{definitionbox}{Definition (Molecular Formula)}
\begin{definition}[Molecular Formula]
\label{def:molecular}
A \emph{molecular formula} is a well-formed formula that is not atomic.
Equivalently, it is a well-formed formula whose final formation step uses a
logical connective or a quantifier.
\end{definition}
\end{definitionbox}

\begin{remark*}[Standard quantified statement]
\[
\varphi\in\mathsf{Form}_{\mathcal L}\setminus\mathsf{Atom}_{\mathcal L}.
\]
\end{remark*}

\begin{remark*}[Definition predicate reading]
\[
\operatorname{MolecularFormula}(\varphi)
\coloneqq
\operatorname{WellFormedFormula}(\varphi)
\land
\neg\operatorname{AtomicFormula}(\varphi).
\]
\end{remark*}

\begin{remark*}[Negated quantified statement]
\[
\varphi\notin\mathsf{Form}_{\mathcal L}
\quad\text{or}\quad
\varphi\in\mathsf{Atom}_{\mathcal L}.
\]
\end{remark*}

\begin{remark*}[Negation predicate reading]
\(\neg\operatorname{MolecularFormula}(\varphi)\) means that \(\varphi\) is
either not a well-formed formula or is atomic.
\end{remark*}

\begin{remark*}[Failure modes]
A string fails to be a molecular formula either because it is not a
well-formed formula at all or because it is an atomic formula with no outer
logical construction.
\end{remark*}

\begin{remark*}[Failure mode decomposition]
\[
\underbrace{x}_{\text{not a formula}}
\qquad
\underbrace{P(x)}_{\text{atomic formula}}
\qquad
\underbrace{P(x)\land}_{\text{not well formed}}.
\]
\end{remark*}

\begin{remark*}[Interpretation]
Molecular formulas are the compound formulas of first-order syntax. They are
formed after the atomic stage by applying a connective or quantifier.
\end{remark*}

\begin{remark*}[Examples]
\(\neg P(x)\), \((P(x)\land Q(x))\), and \(\forall x\,P(x)\) are molecular
formulas.
\end{remark*}

\begin{remark*}[Non-Examples]
\(P(x)\) is atomic rather than molecular, and \(f(x)\) is a term rather than a
formula.
\end{remark*}

\begin{dependencies}
\begin{itemize}
  \item \hyperref[def:wff-pred]{Well-Formed Formula}
  \item \hyperref[def:atomic-formula]{Atomic Formula}
\end{itemize}
\end{dependencies}
